\documentclass{article}

\usepackage[dblblindworkshop]{neurips_2026}
\workshoptitle{Foundation Models for Temporal Systems (FMTS): From Forecasting to World Modeling}

\usepackage[utf8]{inputenc}
\usepackage[T1]{fontenc}
\usepackage{hyperref}
\usepackage{url}
\usepackage{booktabs}
\usepackage{amsfonts}
\usepackage{amsmath}
\usepackage{nicefrac}
\usepackage{microtype}
\usepackage{xcolor}
\usepackage{graphicx}
\usepackage{siunitx}

\title{Forecasting Is Not World Modeling: Accuracy, Action Fidelity, and an Evaluation Matrix for a Task-Critical Thermal Process}

\author{Anonymous}

\begin{document}
\maketitle

\begin{abstract}
Task-critical industrial processes demand more from a temporal model than accurate
forecasts. A model promoted to the inner loop of a controller must respond to candidate
actions with the right sign, magnitude, and transient shape, and it must know when it
cannot. We study this gap on a supercritical once-through boiler, using 82 days of
dual-side field data. A deep forecasting model attains 0.7--0.9~\si{\degreeCelsius} MAE at
18 steps yet attenuates valve responses by 25--450$\times$ relative to the unit's energy
balance, effectively disabling the action channel. A physics-structured world model with
the same data trades accuracy (about 3.5~\si{\degreeCelsius}) for an action channel
that responds on scale at one step but inverts sign at 60 steps---an inversion we
localize to one learned feedback term. We report both findings through a discrimination
matrix:
an identifiability gate on same-type synthetic data, per-layer structural ablations with
budget and convergence diagnostics, and an action-blindness and leakage
audit. Finally we locate the physics model's accuracy cost: after proper training it
approaches persistence at the terminal stage (0.16 vs.\ 0.09~\si{\degreeCelsius}
one-step MAE, beating persistence at 18 steps) while the upstream stage fails
38$\times$ worse than persistence, and the earlier monotone regime bias collapses to
$\pm$0.2~\si{\degreeCelsius}---a training artefact, not a parameter-count one. Our
results support verifiability, not accuracy, as the qualifying property of a world
model, and a two-tier architecture in which a learned proposer is bounded by a
certifiable reduced-order model.
\end{abstract}

\section{Introduction}

A temporal model that forecasts well is not automatically a world model. World models
must maintain a state, imagine futures under interventions, and support decisions
\cite{ha2018world, hafner2023dreamerv3, ding2024understanding}. On safety-critical
plants this qualification question precedes accuracy: a facility can tolerate a
conservative heuristic because it is bounded; it cannot tolerate a model that predicts
confidently outside its own validity \cite{ghosh2026rom}. Closed-loop identification
warned decades ago that passively learned dynamics absorb the confounded structure of
the controller that generated the data \cite{forssell1999closedloop}. Recent evidence
that time-series foundation models collapse under feedback \cite{wan2026forecast},
despite strong open-loop accuracy \cite{ansari2024chronos, das2024timesfm}, suggests
the failure is general, not anecdotal.

We examine the question on a supercritical once-through boiler, a task-critical thermal
process: superheated steam temperature must stay inside a narrow band while disturbances
(load ramps, coal variation, soot blowing) and two spray attemperation valves act on
it. We train two lines on the same dual-side field data. The first is a deep forecaster
with multi-step heads; it reaches 0.7--0.9~\si{\degreeCelsius} MAE at 18 steps. The
second is a physics-structured world model (a lumped-enthalpy state-space with
thermodynamic property tables and learnable parametric closures). The forecaster wins on
accuracy; the world model wins on what the forecaster never learned: responding to
actions. This tension is not a property of our two implementations. It follows from what
the two objectives optimize, and we measure it with a protocol designed to be reused:
the discrimination matrix.

Our contributions are (i) an empirical demonstration, on dual-side field data, that
forecast accuracy does not transfer to world-model qualification; (ii) an evaluation
matrix for such qualification (identifiability gate, per-layer ablations, action-blind
and leakage audits); and (iii) a diagnosis of where the physics model's accuracy cost
comes from, with quantified floors and a convergent agenda.

\section{Testbed and Two Model Lines}

\subsection{The process}

The unit (600~MW class, supercritical, once-through) drives steam through a chain of
superheater stages interleaved with two spray attemperators: low-temperature
superheater, first spray, platen superheater, second spray, final superheater. Five
temperatures along this chain are observed at 10~s resolution, together with boundary
channels (steam flow, coal command, separator pressure and temperature, feedwater
temperature, outlet pressure, spray flow) and two valve positions. The valves are
cross-paired: valve A governs side-B temperatures and vice versa. Data span 82 days per
side; the windows used here are entirely supercritical dry operation. Fast temperature
increments are small (0.09--0.45~\si{\degreeCelsius} per step), so the process is
strongly autoregressive; load-driven swings are orders of magnitude larger.

\subsection{Line 1: a deep forecaster}

Line 1 is a direct multi-step forecaster (sequence model over a history window with
dedicated heads per horizon step, trained with a Gaussian likelihood). It is
representative of the pure forecasting line rather than a state of the art: the horizon
heads are teacher-forced, not recursive. We state this plainly because the comparison
that follows is about what the forecasting objective rewards, and the literature
suggests recursive foundation forecasters inherit the same closed-loop failure
\cite{wan2026forecast}. Its 18-step MAE on held-out windows is 0.712 (side A) and 0.927
(side B).

\subsection{Line 2: a physics-structured world model}

Line 2 embeds the process structure: three lumped-enthalpy segments whose dynamics obey
energy conservation with IAPWS-97 property tables, two spray injections with transport
delays, metal-temperature coupling, and learnable parameters (segment conductances,
evaporation gains, spray efficiencies), in the universal-differential-equation
tradition of embedding structure while learning parameters
\cite{chen2018neuralode, rackauckas2020ude}. An action-blind closure absorbs unmeasured
disturbances from state and whitelisted boundary channels only. A probabilistic
observer estimates the initial state. The model rolls out recursively, so any
counterfactual branch is a consistent trajectory of the same dynamics. Its 18-step MAE
is about 3.5~\si{\degreeCelsius} (v0.1 protocol); retrained values under the v0.2
budget are reported in Section 4.

\section{The Tension, Measured Three Ways}

\subsection{Action response: one line is functionally blind}

We probe both models with $\pm$2\% valve steps, holding all else fixed. Line 1 responds
with $-0.005$ to $-0.015$~\si{\degreeCelsius}, against an energy-balance reference of
roughly $-0.45$ to $-0.87$~\si{\degreeCelsius} for the same 2\%: an attenuation of
25--450$\times$ (Fig.~\ref{fig:evidence}a). Line 2's one-step response on the sprayed
channel follows the energy balance ($-0.2$ to $-0.45$~\si{\degreeCelsius} at 2\%).
Structural honesty requires the rest of the story. The response is partly a
construction guarantee: the action enters an instantaneous mixing equation that the
output map re-evaluates at every step, so the model must respond at step one whether
or not the learned dynamics agree. And at 60 steps the direction breaks: only 12--34\%
of windows end negative (R1's direction rule requires 100\%), with a mean terminal
delta of $+0.09$ to $+0.26$~\si{\degreeCelsius}. An ablation localizes the inversion
to one named term: the wall-rewetting feedback. More spray raises the droplet mass,
which lowers the dry-out factor, which raises the rewetting heat flux into the steam;
the learned rewetting gain over-compensates the direct cooling, and terminal
temperature rises. Zeroing the rewetting coefficient restores the correct sign on
94\% of windows (mean $-0.18$~\si{\degreeCelsius}). The closure is not responsible
(identical probes with and without it). Line 2's action channel therefore exists by
structure but is not yet faithful by physics: a controller using Line 1 as its
internal model sees a plant that does not react to its valves; one using Line 2 sees a
plant that reacts immediately and wrongly at long horizons. Both lines fail the same
requirement---no mechanism certifies the action's sign---through different paths.

\subsection{Counterfactual coherence}

Beyond steady responses, an internal model must carry physically coherent counterfactual
branches. Line 1's teacher-forced heads cannot propagate a perturbation through time;
Line 2's recursive state does so by construction, and its boundary dependence is
whitelisted so that no action leaks into the disturbance path (Section 4, R1). The
confounding is measurable in principle with contemporaneous causal discovery on the
closed-loop data \cite{runge2020pcmci}; our probes below are the model-side equivalent.

\subsection{Accuracy: the cost is real, located, and mostly structural}

Line 2 pays a 3--4$\times$ aggregate accuracy price (H18 MAE
3.9~\si{\degreeCelsius} vs.\ 0.7--0.9 for Line 1), but the price is not uniform
across the steam path. At the terminal stage the retrained model is nearly as good
as persistence one step ahead (0.16 vs.\ 0.09~\si{\degreeCelsius} MAE) and beats it
at 18 steps (1.05 vs.\ 1.18); at the upstream stage it is 38$\times$ worse than
persistence (9.4 vs.\ 0.25~\si{\degreeCelsius}), and the model's learned noise scale
confirms it has written this channel off ($\sigma \approx 8.8$~\si{\degreeCelsius}).
Binning residuals by load shows that the strong monotone regime bias present under
the initial training budget (+2.1 to +5.4~\si{\degreeCelsius} across load bins at the
terminal stage) collapses to $\pm$0.2~\si{\degreeCelsius} after retraining with the
amended budget: level calibration was a training-sufficiency artefact, not a
parameter-count artefact. What remains is structural: the upstream channel's
residual swings $\pm$8.8~\si{\degreeCelsius} across load bins with no monotone
signature, pointing at the enthalpy--temperature inversion at the first segment
(pressure partition assumption) and the anchor coverage (two spray-outlet
temperatures float at initialization), and the first spray-outlet channel retains
44\% regime-correlated variance, the signature of unresolved spray transients
(Fig.~\ref{fig:evidence}). The autoregressive diagnosis generalizes: the process
creeps (0.09--0.45~\si{\degreeCelsius} per step), so large one-step errors are model
jumps, not noise.

\section{A Discrimination Matrix for World-Model Qualification}

The tension above is only meaningful if the measurements are reproducible. We froze an
evaluation protocol, run end-to-end by a single runner against a versioned data
contract, with verdicts computed by paired bootstrap confidence intervals and a 2/3
seed rule.

\textbf{D-SYN identifiability gate.} The full stack (observer, transition, closure)
trains on synthetic data generated by the same architecture and must recover the
teacher: 3/3 seeds pass with NLL reductions of 99--136, well above the 30\% threshold.
Failures on real data therefore reflect structure or data, not an unidentifiable task.

\textbf{O1 initial-state posterior.} A learned posterior must beat the steady-state
anchor by at least 5\% with confidence, and not degrade continuity. Side A verdict:
the learned posterior is MIXED (two seeds improve by 30--31\% with confidence, one
degrades by 13\%), and the precision-weighted hybrid is REJECTED (no seed reaches a
significant improvement: $-0.3\%$, $-1.6\%$, $+3.7\%$, all CIs crossing zero). Under
the v0.2 training budget the steady anchor itself improved enough to erase the
posterior's advantage; the learned posterior's seed instability remains.

\textbf{T1 transition structure.} Nested arms add closure layers and a latent block
against the previous layer. Side A verdicts: the conservative closure is SUPPORTED
(2/3 seeds improve val NLL by 5.7\% and 6.9\% with confidence; seed 3 by 2.6\%,
CI crossing zero). The steam-only placement is REJECTED against it ($-6.4\%$,
$-2.3\%$, $-3.7\%$), and the latent block is REJECTED ($-1.6\%$, $+0.3\%$, $+3.1\%$,
all CIs crossing zero). With the v0.2 budget the latent arm converged to the same
valley the best seed found before, but not consistently enough to earn its
parameters.

\textbf{B1 boundary forecast.} The learned boundary forecaster is benchmarked against
persistence on CRPS. Side A verdict: REJECTED---all three seeds are significantly
worse than persistence (CRPS gaps of 1.7--2.0, CIs entirely below zero). Boundary
channels are strongly autoregressive; a learned forecaster needs more than a GRU to
earn its complexity here.

\textbf{J1 training scheme.} Joint training of transition and observer against staged
training. Side A verdict: SUPPORTED 3/3 (val NLL improvements of 13.5\%, 16.4\%,
33.3\%, all CIs above zero).

\textbf{R1 action-blind and leakage audit.} The closure must not respond to valve
channels (runtime blindness check: 3/3 pass), the boundary whitelist excludes the
spray-flow channel, and leakage probes must not read future information (3/3 clean;
action-aware probes do not improve, so no leakage is suspected). The direction probe
fails: on a 60-step $+5\%$ valve step only 19--34\% of windows end negative (rule
requires 100\%), with mean terminal deltas of $+0.20$ to $+0.26$~\si{\degreeCelsius}.
Verdict: REJECTED on direction; blind and leakage clean.

Each verdict is computed by the frozen runner, not by hand; the matrix doubles as the
artifact we propose as a P5-style honest benchmark for world-model claims
\cite{ghosh2026rom}. All verdicts above are side A; side B is pending at submission
time and will be reported in the camera-ready version.

\begin{figure}[t]
\centering
\includegraphics[width=0.49\textwidth]{figs/fig1_tension.pdf}
\caption{The accuracy--action-fidelity plane. The deep forecaster is accurate but its
valve response is attenuated 25--450$\times$; the physics world model responds on
scale but pays 3--4$\times$ accuracy and, at 60 steps, its terminal response
sign-inverts on most windows (R1).}
\label{fig:tension}
\end{figure}

\begin{figure}[t]
\centering
\includegraphics[width=\textwidth]{figs/fig2_system_matrix.pdf}
\caption{Left: the steam path with five observed temperatures and two cross-paired
spray valves. Right: the discrimination matrix with side-A verdicts.}
\label{fig:system}
\end{figure}

\begin{figure}[t]
\centering
\includegraphics[width=\textwidth]{figs/fig3_evidence.pdf}
\caption{(a) Valve-step responses: the physics line responds on the energy-balance
scale; the deep forecaster is flat. (b) One-step error by channel: the physics line
is near persistence at the terminal stage and 38$\times$ worse upstream. (c) Residual
means by load quintile: the terminal-stage bias is flat after retraining, the
upstream stage swings $\pm$8.8~\si{\degreeCelsius}.}
\label{fig:evidence}
\end{figure}

\section{Discussion}

The retrained physics line already reaches near-persistence quality on the terminal
channel; the remaining cost is concentrated and diagnosable. The upstream channel's
one-step error points at the initialization and inversion layer, not the dynamics:
anchoring all five observations at initialization (two spray-outlet temperatures
currently float) and replacing the fixed pressure-partition assumption with a learned
or measured one targets most of the 9.4~\si{\degreeCelsius} directly. The action-sign
failure and the first spray-outlet channel's residual share one lever: the spray
pathway needs transport and evaporation lag in the mixing term, and the rewetting
feedback needs a sign-and-magnitude contract (its learned gain over-compensates the
direct cooling). Regime-dependent parameter maps, the tempting first remedy, lost
their primary evidence once the training budget was corrected---the monotone load
bias they were meant to absorb was a training artefact. This ordering illustrates the
value of the matrix: one convergence diagnostic and one direction probe re-prioritize
the entire optimization agenda.

For deployment we endorse the convergent architecture the reduced-order-model lineage
articulates \cite{ghosh2026rom, benner2015survey, sha2001simplicity}: a learned
proposer bounded by a certifiable reduced-order model that approves, bounds, or rejects
each proposed action within its certified envelope. Planning with imperfect models is a
documented failure mode \cite{jafferjee2021hallucinating, wang2022doublecheck}; the
two-tier architecture is the constructive answer. The discrimination matrix is the
qualification gate for that proposer. Forecasting remains necessary; it is not
sufficient.

\textbf{Limitations.} Line 1 is teacher-forced, not a recursive foundation forecaster;
our claim about recursive forecasters rests on external evidence
\cite{wan2026forecast} and should be re-measured with one. Data are from a single unit
in dry supercritical operation; two-phase regimes are unobserved. The action-response
probe covers valve steps at fixed load; larger interventions and load-coordinated
moves deserve their own protocol. Matrix verdicts above are single-side (A).

\bibliographystyle{plain}
\bibliography{fmts_refs}

\end{document}
